\section{Conclusion}\label{sec:conclusion}

We presented a two-stage cascaded pipeline for automated Equatorial Plasma Bubble (EPB) detection in $630.0\text{ nm}$ all-sky airglow imagery. By placing an astronomical sky quality control filter (Stage~1, based on high-pass star-pixel counts and mean sky luminance inside the circular dome) ahead of Gaussian background subtraction, CLAHE contrast enhancement, 2DPCA matrix projection ($k=5$, $320$ features), and a heterogeneous stacking ensemble (Stage~2), the system addresses both cloud contamination and low-contrast plasma depletion recognition at a CPU inference cost of $0.63\text{ ms}$ per frame.

Experiments across $572$ labeled exposures and $986$ continuous nocturnal monitoring frames yielded four principal findings. First, the classifier reached $97.20\% \pm 1.57\%$ accuracy ($96.78\% \pm 1.86\%$ F1-score, $0.9984$ ROC-AUC) under 5-fold stratified cross-validation, $95.08\% \pm 5.72\%$ accuracy ($94.24\% \pm 5.37\%$ F1-score, $0.9969$ ROC-AUC) under 5-fold temporal-block group cross-validation across $77$ non-overlapping $30\text{-minute}$ intervals, and $95.80\% \pm 1.57\%$ accuracy when peripheral text stamps were masked out. Second, cross-station transfer experiments showed that while classifiers trained on raw images or flattened pixels achieve high apparent accuracy on random splits by memorizing site-specific luminance and geometry, they fail under zero-shot transfer to an unseen observatory ($0.0\%$ recall for raw images and $54.8\%\text{--}69.0\%$ for flattened RF and 1D-PCA LogReg on Cachoeira Paulista), whereas the preprocessed 2DPCA pipeline achieved $100.0\%$ zero-shot recall ($42/42$). Third, the pipeline maintained $95.88\%$ specificity ($93/97$) under bright lunar glare ($\ge 50\%$ moon illumination), and on an unseen overcast night Stage~1 flagged $79.0\%$ of exposures as cloudy ($90.3\%$ as cloudy or degraded), cutting cloud-induced false EPB triggers by $78.7\%$ to $90.2\%$. Fourth, unsupervised 2DPCA reconstruction error maps and supervised pixel-space SHAP back-projections confirmed that ensemble predictions track physical North-South depletion structures rather than border artifacts.

Future work will expand negative-class sampling at low-latitude margin stations such as Cachoeira Paulista, incorporate multi-year observations across different phases of the solar cycle, and combine optical all-sky classifications with collocated GNSS $S_4$ scintillation receivers and digisonde profiles for real-time space weather forecasting.
